%% =====================================================================
\section{KLN、Bloch--Nordsieck 与红外安全}\label{sec:ir-cancellation}

单个割图的钳夹面给出红外与共线发散；物理截面必须使之相消（对不可分辨的软光子或简并态求和），或把它们吸收进可测量的长程因子。

\begin{theorem}[QED 中的 Bloch--Nordsieck 定理 \tagR]\label{thm:ir-bn-qed}
（$R_{\mathrm{PT}}$：在重整化微扰论中对红外正规化的部分子态严格。）在带质量带电费米子的 QED 中（光子无质量，红外以光子质量 $\mu$ 或维数正规化正规化），对任何外部带电粒子固定的过程，在 $\alpha$ 的每一阶：(a)~对总能量低于分辨率 $E$（在固定参考系中）的实光子求和的截面 $\dd\sigma(E)$ 在去掉红外正规化时有限——虚修正中的软光子发散与实发射中的相消；(b)~依赖于红外正规化的软光子因子在所有阶指数化（Yennie--Frautschi--Suura）。（本定理只涉及微扰 QED。）
% uses: none
\end{theorem}
\begin{proofsketch}
(1)~软光子 $k\to0$ 与外线 $n$（电荷 $eQ_n$，$\eta_n=+1$ 出射、$-1$ 入射）耦合时，振幅因子化为程函因子乘以无光子的振幅，$\mathcal M_{1\gamma}\simeq\mathcal M_0\,e\sum_n\eta_nQ_n\frac{p_n\cdot\varepsilon^*}{p_n\cdot k}$（整体相位依约定而定）；对极化以 $-g_{\mu\nu}$ 求和（程函流 $\sum_n\eta_nQ_np_n^\mu/(p_n\cdot k)$ 由电荷守恒 $\sum_n\eta_nQ_n=0$ 而守恒，故允许）并乘以相空间 $\frac{\dd^3k}{(2\pi)^32\omega}=\frac{\omega\,\dd\omega\,\dd\Omega}{16\pi^3}$，单个软光子的发射概率为 $A\,\dd\omega/\omega$，
\begin{equation}\label{eq:ir-bn-A}
A=-\frac{e^2}{16\pi^3}\int\dd\Omega\sum_{n,m}\eta_n\eta_m\,Q_nQ_m\frac{p_n\cdot p_m}{(p_n\cdot\hat n)(p_m\cdot\hat n)},\qquad \hat n=(1,\hat{\mathbf n});
\end{equation}
对一入一出、$Q=1$、质量为 $m$ 的两个粒子，$A=\frac{2\alpha}{\pi}\big[\frac{1}{2\beta}\ln\frac{1+\beta}{1-\beta}-1\big]$，$\beta=[1-m^4/(p\cdot p')^2]^{1/2}$ 为相对速度，高能时 $A\to\frac{2\alpha}{\pi}[\ln(-q^2/m^2)-1]$（$q=p'-p$）。
(2)~单圈时，虚修正的红外部分为 $\dd\sigma_V=\dd\sigma_0[1-A\ln(\Lambda/\mu)+\dots]$，其中 $\Lambda$ 为软近似失效的能标；能量 $\mu\ll\omega<E$ 的实发射给出 $\dd\sigma_R=\dd\sigma_0[A\ln(E/\mu)+\dots]$；二者之和 $\dd\sigma_0[1+A\ln(E/\Lambda)+\dots]$ 与 $\mu$ 无关。同一个 $A$ 也可写为 $\frac{2\alpha}{\pi}f_{\rm IR}(q^2)$，$f_{\rm IR}$ 即顶点函数红外发散的系数\cite[第~6.4--6.5~节]{PeskinSchroeder1995}。
(3)~高阶时，软光子的虚交换与实发射分别因子化为单圈因子的指数，红外正规化的依赖在两者之积中逐阶相消；对多个实光子能量之和的约束只改变不依赖于 $\mu$ 的因子\cite{BlochNordsieck1937,YennieFrautschiSuura1961,Weinberg1965}。
\cref{eq:ir-bn-A}的角积分与闭式及 $f_{\rm IR}$ 的表示已在三组运动学下数值核对（相对误差 $<10^{-10}$）。
\end{proofsketch}

QED 的相消依赖于软光子耦合只依赖外线的电荷与动量、且光子之间没有相互作用；非阿贝尔理论中后一点不再成立。

\begin{theorem}[KLN 定理 \tagA]\label{thm:ir-kln}
在质量趋于零（或红外正规化去掉）时，对初态\emph{与}末态的简并态集合同时求和的跃迁概率在微扰论的每一阶都是红外有限的。（本定理只涉及部分子态。）
\uses{prop:prelim-unitarity,thm:ir-bn-qed}
\end{theorem}
\begin{proofsketch}
(1)~在含时（老式，old-fashioned）微扰论中，质量奇异性（mass singularity）来自能量分母 $E-E_n$ 在中间态与初态或末态简并时趋于零；简并的含义是：共线的无质量粒子组与单个粒子能量相同，零能粒子的加入不改变能量。
(2)~Kinoshita\cite{Kinoshita1962}与 Lee--Nauenberg\cite{LeeNauenberg1964}论证了：把 $|S_{fi}|^2$ 对初态与末态的全部简并子空间求和后，这些奇异分母逐阶相消；其根源是 $S$ 矩阵在简并子空间上的幺正性（对割图即\cref{prop:prelim-unitarity}的物理幺正性）。对 Yang--Mills 理论的推广见\cite{Sterman:1976jh}，QED 的软光子部分即\cref{thm:ir-bn-qed}。
(3)~$e^+e^-\to$ 强子中初态不带色，只需对末态求和。例如单圈时，以同一归一化 $N_\varepsilon$ 写出，$\sigma(q\bar q)^{(1)}/\sigma_0=N_\varepsilon[-\frac2{\varepsilon^2}-\frac3\varepsilon-8+\pi^2]$，$\sigma(q\bar qg)/\sigma_0=N_\varepsilon[\frac2{\varepsilon^2}+\frac3\varepsilon+\frac{19}2-\pi^2]$，二者之和 $\frac32N_\varepsilon$ 有限，给出 $R=R_0(1+\frac{3C_F}4\frac{\as}\pi)$；而夸克能谱 $\frac1{\sigma_0}\frac{\dd\sigma}{\dd x}$ 保留共线极点 $-\frac1\varepsilon\frac{\as}{2\pi}P_{qq}(x)$——它没有对不同 $x$ 的简并共线态求和（第~6~章）%
% VF-EE-11, VF-EE-6
。
\end{proofsketch}
% VF-EE-11, VF-EE-6

\paragraph{地位与缺口}
一般情形的证明是全阶论证：老式微扰论中简并分母的组合相消在所有阶写出，但未达到 BPHZ 式的完整性。只对一侧（初态或末态）求和也足以相消，但一般需要计入前向散射与不连通的贡献\cite{Frye:2018xjj}。更重要的是物理上的限制：%
% (misconception M8)
强子对撞中初态共线奇异性并不相消，而是被因子化进部分子分布；带色初态的“对初态简并态求和”对强子束流没有物理意义，Bloch--Nordsieck 型的相消在有色初态时甚至会失效（\cref{prop:ir-bn-failure}）。所以 KLN 定理不能解释强子对撞截面为何有限——那是因子化定理的任务（第~12~章）。升级途径：在一般规范理论中把简并态求和的相消写成严格的逐阶陈述。

\begin{proposition}[有色初态时非阿贝尔 Bloch--Nordsieck 抵消的失效 \tagV]\label{prop:ir-bn-failure}
初态为两个有质量的有色夸克（对色平均）时，在两圈阶，两个软胶子产生的次领头红外发散在 Bloch--Nordsieck 意义下的单举截面中不相消；不相消的项在高能下是幂次压低的。（本命题只涉及部分子态。）
\uses{thm:ir-bn-qed,def:prelim-feynman-rules}
\end{proposition}
\begin{proofsketch}
定性地说，其结构如下（这是文献中对该机制的通常解释）。在阿贝尔理论中，两个入射带电粒子之间软光子交换的虚部（Coulomb 相位）是一个纯相位，在 $|\mathcal M|^2$ 中消失，其余的程函因子与实发射相消（\cref{thm:ir-bn-qed}）。在非阿贝尔理论中，Coulomb（Glauber）相位乘的是两入射夸克之间的色矩阵 $\mathbf T_1\cdot\mathbf T_2$；两圈时它与色关联的软实发射相干涉，而这些色矩阵彼此不对易，留下即使对初态色平均之后也不相消的红外发散项，它们随质量趋于零（固定 $s$）而压低。显式的两圈计算见\cite{DoriaFrenkelTaylor1980,DiLietoEtAl1981,Andrasi:1980qw}。
\end{proofsketch}

\paragraph{地位与缺口}
证据类型：显式的微扰计算（两圈，有质量有色初态），即 \tagV 的依据是一个显式的 QCD 反例。据其摘要，Doria、Frenkel 与 Taylor 发现两个有质量、色平均的初态夸克中两个软胶子产生的次领头红外发散不相消\cite{DoriaFrenkelTaylor1980}，Di'Lieto 等人确认了不相消（但中间结果有所不同）\cite{DiLietoEtAl1981}，Andrási 等人亦然\cite{Andrasi:1980qw}。据其摘要，不相消的项在高能下幂次压低：Frenkel 等人在 $O(g^4)$ 证明了这一点并论证了全阶推广\cite{Frenkel:1983di}，Catani 等人指出它们一般是高扭度型的\cite{Catani:1985xt}；CSS 的综述写道，在 Drell--Yan 过程中两圈的不相消发散出现在 $m^4/s^2$ 的层次，而微扰论在 $m^2/s$ 的层次尊重因子化\cite[第~10.2~节]{CollinsSoperSterman1989}。这一失效的前提与 Drell--Yan 因子化定理的前提不同——前者是固定色、有质量的有色入射部分子，后者是色单态强子中的无质量部分子——见第~12~章的\cref{rem:dy-dft-premises}。因此本命题并不与领头幂次的 Drell--Yan 因子化（全阶论证）矛盾，但它说明非阿贝尔理论中软胶子的相消不能简单搬用 QED 的论证。

\begin{definition}[红外安全可观测量]\label{def:ir-safe}
$e^+e^-\to$ 强子中的一个可观测量由一列对称函数 $\mathcal S_n(p_1,\dots,p_n)$（$n=2,3,\dots$，$p_i$ 为末态无质量部分子的动量）给出，其加权截面为 $\sigma[\mathcal S]=\sum_n\int\dd\Phi_n\,|\mathcal M_n|^2\,\mathcal S_n$。称它是\emph{红外安全}的，若每个 $\mathcal S_n$ 连续，并且：(i)~（软）对任意方向，$\mathcal S_{n+1}(p_1,\dots,p_n,\epsilon p)\to\mathcal S_n(p_1,\dots,p_n)$（$\epsilon\to0$）；(ii)~（共线）$\mathcal S_{n+1}(p_1,\dots,p_{n-1},zp,(1-z)p)=\mathcal S_n(p_1,\dots,p_{n-1},p)$（$0\le z\le1$）。这是红外安全性（infrared safety）的非形式定义\cref{def:qcd-ir-safety-informal}的精确形式：权重只依赖于能流，对零能部分子的加入与严格共线的分裂都不敏感。例：总截面（$\mathcal S_n\equiv1$）、推力等事件形状、Sterman--Weinberg 锥形喷注截面（第~16~章）；反例：部分子数、单个夸克的能谱。
\uses{def:qcd-ir-safety-informal}
\end{definition}

\begin{theorem}[红外安全可观测量的有限性 \tagA]\label{thm:ir-safe-finite}
红外安全的 $e^+e^-$ 可观测量在无质量微扰论中逐阶有限。（本定理只涉及部分子态。）
% errtype: E-S
\uses{def:ir-safe,thm:ir-leading-regions,prop:prelim-unitarity}
\end{theorem}
\begin{proofsketch}
(1)~加权截面是割开的真空极化图对全部割线求和、每个割线乘以相应末态权重的结果（\cref{prop:prelim-unitarity}）。
(2)~每个未割图的钳夹面由\cref{thm:ir-leading-regions}的约化图给出：硬顶点、末态喷注与软子图，且在超领头项相消后至多对数发散。
(3)~在一个钳夹面附近，同一未割图的不同割法（切过喷注内部或软子图的不同线）对应同一能流：喷注中的共线部分子的总动量、软部分子的零能量。由红外安全性，这些割法的权重在钳夹面上相等；于是对割线的求和在权重相等的部分中重现未割图的幺正性关系，钳夹面上的奇异性在求和中相消，剩下的部分正比于权重之差，在钳夹面上为零，从而可积。
文献：\cite{Sterman1978a,Sterman1978b,Sterman:1979uw,StermanWeinberg1977,Ore:1979ry}。
\end{proofsketch}

\paragraph{地位与缺口}
该论证依赖领头区域的分析（\cref{thm:ir-leading-regions}，\tagA）以及割线求和中末态奇异性的相消，后者只在 E-S 层次上给出（权重差在钳夹面附近的消失速度与幂次计数的结合未写成积分界）\cite{Sterman1978a,Sterman1978b}。总截面另有独立的严格证明（第~6~章的\cref{thm:ee-R-finite}，\tagR）。红外安全只处理末态奇异性：%
% (misconception M9, first statement)
初态强子仍需要因子化，而红外安全的观测量也可能有非全局对数、超领头对数、Glauber 敏感性与 $O(\Lambda/Q)$ 的强子化修正（第~16、21~章）。升级途径：把 (3) 写成对嵌套邻域一致的积分界。
